Este capítulo detalla el proceso de construcción del sistema \textbf{DOMU}, describiendo la estructura técnica del código fuente, la configuración del entorno, las fases de implementación ejecutadas bajo la metodología híbrida definida y los aspectos técnicos más relevantes de la solución. Se aborda la evolución del software desde su núcleo de autenticación hasta los módulos de gestión comunitaria, financiera y de comunicación en tiempo real.

\section{Estructura del Código y Arquitectura}
\label{sec:estructura_codigo}

El sistema se ha desarrollado siguiendo una arquitectura cliente-servidor desacoplada. El Backend expone una API REST con más de 120 endpoints, consumida por el Frontend construido como una \textit{Single Page Application} (SPA). Ambos repositorios se gestionan de forma independiente mediante Git, lo que permite ciclos de desarrollo, pruebas y despliegue autónomos.

\subsection{Backend (Java)}
El servidor se ha construido con \textbf{Java~21} y el framework ligero \textbf{Javalin~6}~\cite{javalin2025}, gestionado con \textbf{Gradle} (Kotlin DSL). El proyecto cuenta con aproximadamente 190~clases Java organizadas en una arquitectura por capas. La estructura principal bajo el paquete raíz \texttt{com.domu} es la siguiente:

\begin{itemize}
    \item \textbf{config/}: Configuración global del sistema. La clase \texttt{AppConfig} centraliza los parámetros del servidor, base de datos, JWT, correo electrónico y almacenamiento en Box. La clase \texttt{DependencyInjectionModule} configura el contenedor de inyección de dependencias mediante \textbf{Google Guice}, registrando como \textit{singletons} los 27~servicios, 21~repositorios, los manejadores WebSocket y el módulo de seguridad.

    \item \textbf{domain/}: Contiene 35~entidades de dominio organizadas en subpaquetes temáticos que reflejan los módulos del sistema:
    \begin{itemize}
        \item \texttt{core/}: Entidades base (\texttt{User}, \texttt{Building}, \texttt{HousingUnit}, \texttt{Role}).
        \item \texttt{access/}: Control de acceso (\texttt{Visit}, \texttt{VisitAuthorization}, \texttt{AccessLog}, \texttt{DeliveryPackage}).
        \item \texttt{finance/}: Gestión financiera (\texttt{CommonExpensePeriod}, \texttt{CommonCharge}, \texttt{CommonPayment}, \texttt{DelinquencyRecord}).
        \item \texttt{facility/}: Áreas comunes (\texttt{CommonArea}, \texttt{Reservation}).
        \item \texttt{voting/}: Votaciones (\texttt{VotingEvent}, \texttt{VotingOption}, \texttt{Vote}).
        \item \texttt{community/}: Foro (\texttt{ForumCategory}, \texttt{ForumThread}, \texttt{ForumMessage}).
        \item \texttt{staff/}: Personal (\texttt{Personnel}, \texttt{Shift}, \texttt{Task}, \texttt{PersonnelLog}).
        \item \texttt{ticket/}: Incidentes (\texttt{Ticket}, \texttt{TicketFollowUp}).
        \item \texttt{vendor/}: Proveedores (\texttt{Provider}, \texttt{ServiceRequest}, \texttt{Quote}, \texttt{Invoice}, \texttt{PreventiveMaintenance}, \texttt{FinancialTransaction}).
    \end{itemize}

    \item \textbf{dto/}: Aproximadamente 99~\textit{Data Transfer Objects} que definen los contratos de la API. Cada módulo cuenta con DTOs de solicitud (\texttt{*Request}) y respuesta (\texttt{*Response}), separando la representación HTTP de las entidades de dominio. Ejemplos representativos incluyen \texttt{LoginRequest}, \texttt{AuthResponse}, \texttt{CreateVisitRequest}, \texttt{CommonExpensePeriodDetailResponse} y \texttt{ChatMessageResponse}.

    \item \textbf{database/}: Capa de persistencia compuesta por 21~repositorios que implementan el acceso a datos mediante JDBC directo sobre MySQL, utilizando \textbf{HikariCP}~5.1 para el \textit{pool} de conexiones. Las clases principales incluyen \texttt{UserRepository}, \texttt{CommonExpenseRepository}, \texttt{VisitRepository}, \texttt{PollRepository}, \texttt{AmenityRepository}, \texttt{ChatRepository}, \texttt{ParcelRepository}, \texttt{NotificationRepository}, entre otros. La clase \texttt{DataSourceFactory} centraliza la creación del \textit{pool} de conexiones.

    \item \textbf{service/}: Capa de lógica de negocio con 27~servicios que encapsulan las reglas del sistema. Los servicios más relevantes son:
    \begin{itemize}
        \item \texttt{UserService}: Registro, autenticación, confirmación por correo y recuperación de contraseñas.
        \item \texttt{BuildingService}: Flujo de registro de comunidades con aprobación por correo e invitación al administrador.
        \item \texttt{CommonExpenseService}: Creación de períodos de gasto común, cargos por unidad y procesamiento de pagos.
        \item \texttt{CommonExpensePdfService}, \texttt{PaymentReceiptPdfService}, \texttt{ChargeReceiptPdfService}: Generación de documentos PDF con \textbf{OpenPDF}.
        \item \texttt{VisitService}: Autorización de visitas, generación de códigos QR y \textit{check-in}.
        \item \texttt{ChatService}: Gestión de salas de chat y mensajes en tiempo real.
        \item \texttt{NotificationService}: Creación y distribución de notificaciones vía WebSocket.
    \end{itemize}

    \item \textbf{security/}: Módulo de seguridad compuesto por 4~clases: \texttt{JwtProvider} (generación y validación de tokens JWT con HMAC256), \texttt{BCryptPasswordHasher} (hashing de contraseñas), \texttt{PasswordHasher} (interfaz de abstracción) y \texttt{AuthenticationHandler} (\textit{middleware} de Javalin que intercepta las rutas protegidas y valida el token en la cabecera \texttt{Authorization}).

    \item \textbf{email/}: Servicio de correo electrónico con 6~clases: la interfaz \texttt{EmailService}, su implementación SMTP (\texttt{SmtpEmailService}) y cuatro plantillas HTML para confirmación de cuenta, aprobación de comunidades, notificación de aprobación y notificación de rechazo.

    \item \textbf{web/}: Capa de presentación HTTP con 4~clases. \texttt{WebServer} es la clase principal que registra los más de 120~endpoints REST y los dos canales WebSocket (\texttt{/ws/chat} y \texttt{/ws/notifications}). \texttt{ChatWebSocketHandler} y \texttt{NotificationWebSocketHandler} gestionan la comunicación bidireccional en tiempo real. \texttt{UserMapper} convierte entidades de usuario a DTOs de respuesta.
\end{itemize}

El esquema de base de datos se versiona mediante \textbf{Flyway}, con 35~migraciones SQL que documentan la evolución del modelo de datos desde la creación del esquema de autenticación (V001) hasta las últimas actualizaciones del sistema de notificaciones y estructura de edificios (V035).

\subsection{Frontend (React)}
La interfaz de usuario se ha construido con \textbf{React~19} y el empaquetador \textbf{Vite~7}, utilizando JavaScript (JSX) y hojas de estilo SCSS. El proyecto comprende aproximadamente 92~archivos fuente con más de 20.000~líneas de código, organizados de la siguiente manera:

\begin{itemize}
    \item \textbf{pages/} (54~vistas): Componentes de nivel de ruta que conforman las pantallas de la aplicación. Están segmentadas por rol:
    \begin{itemize}
        \item \textit{Públicas} (7): \texttt{Home}, \texttt{Login}, \texttt{Register}, \texttt{ForgotPassword}, \texttt{About}, \texttt{Contact}, \texttt{UserTypeProveedores}.
        \item \textit{Administrador} (13): \texttt{Dashboard}, \texttt{AdminCreateUser}, \texttt{AdminIncidentsBoard}, \texttt{AdminIncidentStats}, \texttt{AdminResidents}, \texttt{AdminHousingUnits}, \texttt{AdminAmenities}, \texttt{AdminCommonExpenses}, \texttt{AdminParcels}, \texttt{AdminTasks}, \texttt{AdminStaff}, \texttt{AdminProviders}, \texttt{AdminServiceOrders}.
        \item \textit{Residente} (18): \texttt{ResidentProfile}, \texttt{ResidentVisits}, \texttt{ResidentIncidents}, \texttt{ResidentAmenities}, \texttt{ResidentCartola}, \texttt{ResidentChargesDetail}, \texttt{ResidentPaymentFlow}, \texttt{ResidentParcels}, \texttt{ResidentPublications}, \texttt{ResidentLibrary}, \texttt{ResidentMarketplace}, \texttt{ChatHub}, \texttt{Votaciones}, entre otras.
        \item \textit{Proveedor} (2): \texttt{ProviderPortal}, \texttt{ProviderServiceOrders}.
        \item \textit{Especiales}: Páginas de soluciones por tipo de usuario, confirmación de correo, registro de administrador por invitación y centro de notificaciones.
    \end{itemize}

    \item \textbf{components/} (18~componentes): Elementos reutilizables de la interfaz, incluyendo \texttt{Button} (con estados de carga), \texttt{FormField} (con validación), \texttt{Spinner}, \texttt{Skeleton} (\textit{skeleton loading}), \texttt{ConfirmModal}, \texttt{LocationPicker} (selector de ubicación basado en mapas mediante la librería \texttt{pigeon-maps}), \texttt{VisitRegistrationPanel}, \texttt{MarketCard}, \texttt{ProtectedRoute} y un módulo completo de pago (\texttt{PaymentMethodModal}, \texttt{CardPaymentView}, \texttt{BankTransferView}, \texttt{ChargeSelector}, \texttt{PaymentSuccess}).

    \item \textbf{layout/} (7~componentes): Estructura visual de la aplicación, compuesta por \texttt{ProtectedLayout} (envolvente para rutas autenticadas con barra lateral y cabecera), \texttt{AuthLayout} (envolvente para páginas públicas), \texttt{Header}, \texttt{Sidebar} (navegación filtrada por rol), \texttt{Footer} y \texttt{MainContent}.

    \item \textbf{services/}: Capa de comunicación con el Backend. El archivo principal \texttt{api.js} (1.034~líneas) centraliza todas las llamadas HTTP organizadas por módulo (\texttt{auth}, \texttt{visits}, \texttt{incidents}, \texttt{finance}, \texttt{polls}, \texttt{amenities}, \texttt{chat}, \texttt{market}, \texttt{parcels}, \texttt{library}, \texttt{notifications}, entre otros). Cada petición inyecta automáticamente el token JWT en la cabecera \texttt{Authorization} y el identificador de la comunidad activa en la cabecera \texttt{X-Building-Id}.

    \item \textbf{context/}: Estado global de la aplicación mediante \textbf{React Context API}. El \texttt{AppContext} gestiona la sesión del usuario autenticado, el edificio seleccionado (soporte multi-comunidad), el tema visual (claro/oscuro) y funciones de \textit{logout} y cambio de comunidad.

    \item \textbf{hooks/}: \textit{Custom hooks} de React. El hook \texttt{useNotifications} integra WebSocket con \textbf{TanStack React Query}~v5 para recibir notificaciones en tiempo real, con reconexión automática cada 5~segundos ante desconexiones.

    \item \textbf{constants/}: Configuración centralizada con 85~definiciones de rutas, secciones de navegación filtradas por rol y 20~tipos de notificación con iconos y colores asociados.

    \item \textbf{styles/}: Hojas de estilo globales en SCSS compiladas por Vite.
\end{itemize}

La aplicación se configura como \textit{Progressive Web App} (PWA) mediante el plugin \texttt{vite-plugin-pwa}, y utiliza TanStack React Query para la gestión de caché y revalidación de datos del servidor.


\section{Fases de Implementación (Iteraciones)}
\label{sec:fases_implementacion}

Siguiendo la metodología híbrida adoptada, el desarrollo se dividió en fases iterativas que permitieron desplegar funcionalidades incrementales. Cada fase fue validada antes de avanzar a la siguiente, asegurando la estabilidad del sistema. A continuación, se describen las seis iteraciones ejecutadas.

\subsection{Fase 1: Fundamentos y Seguridad}
El objetivo de esta primera iteración fue establecer los cimientos técnicos del sistema y asegurar el acceso controlado.
\begin{itemize}
    \item \textbf{Configuración del Entorno:} Se estableció el proyecto Backend con Gradle, configurando el servidor Javalin y la conexión a MySQL mediante HikariCP. Se implementó Flyway para el versionamiento de migraciones de base de datos, asegurando la reproducibilidad del esquema entre entornos de desarrollo y producción.
    \item \textbf{Inyección de Dependencias:} Se configuró Google Guice como contenedor de inversión de control, registrando todos los servicios y repositorios como \textit{singletons} en el módulo \texttt{DependencyInjectionModule}.
    \item \textbf{Autenticación (JWT):} Se implementó un sistema de seguridad basado en \textit{JSON Web Tokens} firmados con HMAC256. Al iniciar sesión, el usuario recibe un token con una duración configurable (por defecto 60~minutos), el cual se almacena en el \texttt{localStorage} del navegador. El \textit{middleware} \texttt{AuthenticationHandler} intercepta todas las rutas protegidas y valida la vigencia del token.
    \item \textbf{Registro con Confirmación:} Se implementó el flujo completo de registro de usuarios con confirmación por correo electrónico, incluyendo el servicio SMTP con plantillas HTML (\texttt{UserConfirmationEmailTemplate}) y la recuperación de contraseñas mediante tokens de un solo uso.
    \item \textbf{Gestión de Perfiles:} Funcionalidad para que los usuarios puedan editar su información personal, cambiar su contraseña y subir avatar de perfil, con optimización automática de imágenes mediante la librería Thumbnailator.
\end{itemize}

\subsection{Fase 2: Gestión de Comunidades y Usuarios}
Se desarrollaron las herramientas administrativas necesarias para la configuración y operación de las comunidades.
\begin{itemize}
    \item \textbf{Registro de Comunidades:} Se implementó un flujo completo de solicitud, aprobación y rechazo de comunidades. La solicitud se envía mediante un formulario y la documentación se almacena en la nube mediante \textbf{Box}. El aprobador recibe un correo electrónico con enlaces para aprobar o rechazar la solicitud, tras lo cual se genera automáticamente una invitación para que el administrador de la comunidad registre su cuenta.
    \item \textbf{Soporte Multi-Comunidad:} Se diseñó un mecanismo que permite a un usuario pertenecer a múltiples edificios. El Frontend envía la cabecera \texttt{X-Building-Id} en cada petición, y el Backend filtra los datos según la comunidad activa. El contexto global del Frontend gestiona el cambio de comunidad e invalida las cachés de React Query automáticamente.
    \item \textbf{Gestión de Unidades Habitacionales:} CRUD completo de unidades (departamentos o casas) dentro de cada edificio, con la posibilidad de vincular y desvincular residentes a cada unidad.
    \item \textbf{Creación de Usuarios por Rol:} Interfaces para que el administrador pueda crear y gestionar cuentas de Conserjes, Personal de Aseo y otros funcionarios. El sistema cuenta con 5~roles base: Administrador, Residente, Conserje, Funcionario y Proveedor.
    \item \textbf{Vistas Segmentadas por Rol:} Se implementó la segmentación de la interfaz mostrando secciones de navegación y funcionalidades específicas según el rol del usuario autenticado, tanto en el \texttt{Sidebar} como en las rutas protegidas del Frontend.
\end{itemize}

\subsection{Fase 3: Servicios al Residente y Operación}
Esta fase se centró en las funcionalidades de interacción diaria dentro de la comunidad.
\begin{itemize}
    \item \textbf{Reporte y Gestión de Incidentes:} Módulo que permite a los residentes reportar problemas (filtraciones de agua, ruidos molestos, fallas eléctricas, entre otros), generando un ticket con categoría, prioridad y descripción. La administración visualiza los incidentes en un tablero tipo Kanban (\texttt{AdminIncidentsBoard}) con estados (abierto, en progreso, resuelto) y puede asignarlos a personal específico. Se implementó también una vista de estadísticas (\texttt{AdminIncidentStats}) que presenta métricas de los incidentes.
    \item \textbf{Reservas de Áreas Comunes:} Sistema de calendario que permite al administrador crear áreas comunes (quinchos, salones, gimnasios), configurar franjas horarias (\textit{time slots}) y definir aforos máximos. Los residentes pueden consultar la disponibilidad por fecha y reservar espacios desde la interfaz. Se implementó la lógica de cancelación y la restricción de reserva por morosidad.
    \item \textbf{Sistema de Votaciones:} Módulo que permite a la administración crear eventos de votación con múltiples opciones, establecer fechas de cierre y publicar resultados. Los residentes emiten su voto de manera remota. Se implementó el cierre programado de votaciones y la exportación de resultados a CSV para respaldo y auditoría.
    \item \textbf{Control de Acceso y Visitas:} Se desarrolló el módulo de autorización de visitas, permitiendo a los residentes pre-autorizar visitantes con fecha, hora y motivo. El conserje puede realizar el \textit{check-in} mediante lectura de código QR (endpoint \texttt{POST /api/visits/qr-check}) o de forma manual. Se implementó la gestión de contactos frecuentes (\texttt{VisitContactService}), que permite a los residentes guardar datos de visitantes habituales y generar visitas rápidamente a partir de ellos.
\end{itemize}

\subsection{Fase 4: Módulo Financiero}
Esta fase abordó los módulos de mayor complejidad en la lógica de negocio.
\begin{itemize}
    \item \textbf{Períodos de Gastos Comunes:} Se implementó la lógica para que el administrador cree períodos mensuales de gasto común y agregue cargos individuales por unidad habitacional, considerando el coeficiente de prorrata de cada departamento.
    \item \textbf{Flujo de Pagos:} Desarrollo completo del flujo transaccional que guía al residente a través del proceso de pago. El Frontend presenta un selector de cargos pendientes (\texttt{ChargeSelector}), opciones de método de pago (\texttt{PaymentMethodModal}) y vistas específicas para transferencia bancaria y tarjeta. El Backend procesa el pago y actualiza el estado del cargo.
    \item \textbf{Generación de Documentos PDF:} Mediante la librería \textbf{OpenPDF}, se implementaron tres servicios de generación de documentos: estados de cuenta mensuales (\texttt{CommonExpensePdfService}), comprobantes de pago (\texttt{PaymentReceiptPdfService}) y boletas de cargo (\texttt{ChargeReceiptPdfService}). Los residentes pueden descargar estos documentos directamente desde la interfaz.
    \item \textbf{Subida de Comprobantes:} Los residentes pueden subir comprobantes de transferencia bancaria, los cuales se almacenan en \textbf{Box} y quedan vinculados al cargo correspondiente para su posterior verificación por la administración.
\end{itemize}

\subsection{Fase 5: Comunicación y Módulos Comunitarios}
Se implementaron las funcionalidades orientadas a la interacción entre residentes y la vida comunitaria.
\begin{itemize}
    \item \textbf{Chat en Tiempo Real:} Sistema de mensajería instantánea construido sobre \textbf{WebSocket} (canal \texttt{/ws/chat}). El \texttt{ChatWebSocketHandler} gestiona la conexión bidireccional autenticada mediante token JWT. Los residentes pueden iniciar conversaciones con vecinos, y el sistema incluye un flujo de solicitud de chat para respetar la privacidad. Se implementó la visualización de usuarios en línea y el ocultamiento de salas.
    \item \textbf{Foro Comunitario:} Módulo de publicaciones que permite crear hilos de discusión categorizados, editarlos y eliminarlos. Funciona como un espacio de comunicación asincrónica complementario al chat.
    \item \textbf{Marketplace Comunitario:} Tienda interna que permite a los residentes publicar artículos en venta con múltiples imágenes, descripción y precio. Las imágenes se almacenan en la nube mediante \texttt{MarketplaceStorageService}. Se integra con el chat para facilitar la comunicación entre comprador y vendedor.
    \item \textbf{Gestión de Encomiendas:} Módulo para que el conserje registre la recepción de paquetes, asignándolos a la unidad habitacional correspondiente. El residente recibe una notificación y puede consultar el estado de sus encomiendas. Se implementaron cambios de estado (recibido, notificado, entregado) y la vista administrativa (\texttt{AdminParcels}).
    \item \textbf{Biblioteca de Documentos:} Repositorio digital que permite al administrador subir documentos relevantes para la comunidad (reglamentos, actas, comunicados). Los archivos se almacenan en Box y los residentes pueden descargarlos desde la interfaz.
\end{itemize}

\subsection{Fase 6: Personal, Notificaciones y Cierre}
La última iteración completó los módulos de gestión operativa y el sistema de notificaciones.
\begin{itemize}
    \item \textbf{Gestión de Personal:} CRUD completo para el registro de funcionarios (conserjes, personal de aseo, mayordomos), incluyendo la administración de turnos y la visualización del personal activo.
    \item \textbf{Asignación de Tareas:} Módulo que permite al administrador crear tareas, asignarlas a funcionarios específicos y realizar seguimiento de su estado (pendiente, en progreso, completada). Los funcionarios pueden consultar sus tareas asignadas desde la vista \texttt{StaffTasks}.
    \item \textbf{Sistema de Notificaciones en Tiempo Real:} Se implementó un canal WebSocket dedicado (\texttt{/ws/notifications}) gestionado por \texttt{NotificationWebSocketHandler}. El Backend genera notificaciones ante eventos relevantes (nueva encomienda, pago recibido, incidente actualizado, nueva visita, entre otros). El Frontend las recibe mediante el hook \texttt{useNotifications}, que integra la conexión WebSocket con TanStack React Query para mantener el estado sincronizado. Se implementó también el centro de notificaciones (\texttt{NotificationCenter}) con funcionalidad de marcar como leído y preferencias de notificación por tipo.
    \item \textbf{Servicio de Correo Electrónico:} Se consolidó el sistema de correo con plantillas HTML para confirmación de cuenta, aprobación/rechazo de comunidades y notificación al administrador invitado, utilizando Jakarta Mail sobre SMTP.
\end{itemize}

\subsection{Fase 7: Proveedores, Órdenes de Servicio y Documentación de la API}
Esta fase incorporó el módulo de gestión de proveedores externos y la documentación interactiva de la API.
\begin{itemize}
    \item \textbf{Gestión de Proveedores:} CRUD completo para registrar empresas proveedoras de servicios vinculadas a la comunidad (electricistas, gasfíteres, empresas de aseo, entre otros). El administrador puede crear, editar y desactivar proveedores desde la vista \texttt{AdminProviders}. Cada proveedor se asocia a una categoría de servicio y opcionalmente a una cuenta de usuario con rol \textit{Proveedor}.
    \item \textbf{Órdenes de Servicio:} Flujo completo de solicitud de trabajo que permite al administrador crear órdenes, asignarlas a un proveedor y realizar seguimiento por estado (pendiente, aceptada, rechazada, en progreso, completada, cancelada). El proveedor puede aceptar o rechazar órdenes, subir cotizaciones y marcar trabajos como completados desde su portal (\texttt{ProviderPortal} y \texttt{ProviderServiceOrders}).
    \item \textbf{Cotizaciones:} Los proveedores pueden adjuntar propuestas económicas a cada orden de servicio, incluyendo monto, descripción y archivo adjunto. El administrador revisa y aprueba o rechaza las cotizaciones.
    \item \textbf{Documentación de la API (Swagger~UI):} Se integró el plugin comunitario \texttt{javalin-openapi} para generar automáticamente la especificación OpenAPI~3 del sistema. La interfaz Swagger~UI, accesible en \texttt{/swagger}, permite explorar y probar los más de 130~endpoints de la API de forma interactiva.
\end{itemize}

\subsection{Contratos de la API REST}
Los contratos detallados de los endpoints principales de la API ---incluyendo
métodos HTTP, rutas, objetos de solicitud (\textit{request}) y respuesta
(\textit{response})--- se documentan en el Anexo~\ref{anexo:endpoints}.
Dicha documentación complementa la especificación interactiva disponible
en Swagger~UI y facilita la comprensión de la integración entre el Frontend
y el Backend para cada módulo funcional del sistema.

\section{Detalles Técnicos de Implementación}

\subsection{Integración del Lector de Cédula}
Para el control de acceso en conserjería, se implementó una estrategia de integración mediante hardware tipo ``pistola escáner'' de códigos QR/Barras. Este dispositivo opera en modo de emulación de teclado (\textit{HID Mode}), lo que permite su uso directo sobre el navegador web sin necesidad de controladores especiales.

El flujo técnico implementado es el siguiente:
\begin{enumerate}
    \item El conserje accede a la vista de registro de visitas en la interfaz web (\texttt{VisitRegistrationPanel}).
    \item Al escanear el código QR de la Cédula de Identidad chilena, el dispositivo ``escribe'' automáticamente la cadena de texto contenida en el código.
    \item El Frontend procesa esta cadena, extrayendo el RUN del visitante mediante \textit{parsing} de la estructura del código.
    \item Se realiza una consulta automática al Backend (endpoint \texttt{POST /api/visits/qr-check}) que verifica si la persona tiene una autorización vigente. En caso positivo, se completa el \textit{check-in} automáticamente; en caso negativo, se presenta el formulario para registro manual.
    \item Paralelamente, el sistema de contactos frecuentes (\texttt{POST /api/visit-contacts/\{id\}/register}) permite generar visitas con un solo clic para visitantes previamente guardados.
\end{enumerate}

\subsection{Manejo de Seguridad y Sesión}
La seguridad del sistema se gestiona en múltiples capas:

\begin{itemize}
    \item \textbf{Autenticación:} La clase \texttt{JwtProvider} genera tokens JWT firmados con el algoritmo HMAC256. El token contiene el identificador del usuario, su correo electrónico y una fecha de expiración configurable. Las contraseñas se almacenan encriptadas mediante \texttt{BCrypt} con sal automática.

    \item \textbf{Middleware de Autorización:} La clase \texttt{AuthenticationHandler} actúa como un \textit{before-handler} de Javalin que intercepta las peticiones a rutas protegidas. Extrae el token de la cabecera \texttt{Authorization: Bearer <token>}, lo valida y deposita el identificador del usuario en el contexto de la petición para su uso posterior por los controladores.

    \item \textbf{Interceptor en el Frontend:} El módulo \texttt{services/api.js} inyecta automáticamente el token JWT y la cabecera \texttt{X-Building-Id} en cada petición. Si el servidor responde con un código 401 (no autorizado) debido a la expiración del token, el Frontend cierra la sesión automáticamente, limpia el \texttt{localStorage} y redirige al formulario de inicio de sesión.

    \item \textbf{Recuperación de Contraseña:} Se implementó un flujo seguro de recuperación que genera un token de un solo uso almacenado en la tabla \texttt{password\_reset\_tokens}, el cual se envía por correo electrónico al usuario y expira tras su primer uso.
\end{itemize}

\subsection{Comunicación en Tiempo Real}
El sistema implementa dos canales WebSocket nativos de Javalin para comunicación bidireccional:

\begin{itemize}
    \item \textbf{Chat (\texttt{/ws/chat}):} Gestionado por \texttt{ChatWebSocketHandler}, permite el intercambio de mensajes instantáneos entre residentes. La conexión se autentica mediante el token JWT enviado como parámetro de consulta (\texttt{?token=...}). Los mensajes se persisten en la base de datos a través de \texttt{ChatRepository} y se distribuyen en tiempo real a los participantes de la sala.

    \item \textbf{Notificaciones (\texttt{/ws/notifications}):} Gestionado por \texttt{NotificationWebSocketHandler}, distribuye notificaciones a los usuarios conectados cuando ocurren eventos relevantes en el sistema. En el Frontend, el hook \texttt{useNotifications} mantiene la conexión WebSocket y sincroniza el estado con TanStack React Query, implementando reconexión automática cada 5~segundos ante desconexiones.
\end{itemize}

\subsection{Generación de Documentos PDF}
Para la generación de documentos financieros se utiliza la librería \textbf{OpenPDF}~1.3, que permite construir documentos PDF de forma programática. Se implementaron tres servicios especializados:

\begin{itemize}
    \item \texttt{CommonExpensePdfService}: Genera el estado de cuenta mensual de gastos comunes, detallando los cargos por unidad, el monto total y el período correspondiente.
    \item \texttt{PaymentReceiptPdfService}: Genera comprobantes de pago con los datos del residente, el monto pagado, la fecha y el método de pago.
    \item \texttt{ChargeReceiptPdfService}: Genera boletas individuales por cada cargo registrado.
\end{itemize}

Los documentos se generan bajo demanda cuando el residente los solicita desde la interfaz y se entregan como descarga directa a través de la API.

\subsection{Almacenamiento en la Nube}
El almacenamiento de archivos se gestiona mediante \textbf{Box}, utilizando el SDK oficial de Java con autenticación por \textit{Developer Token}. La clase central \texttt{BoxStorageService} encapsula las operaciones de subida, descarga y organización de archivos en carpetas, manteniendo una caché interna de identificadores de carpetas (\texttt{folderIdCache}) para evitar consultas redundantes a la API de Box.

Tres servicios especializados delegan en \texttt{BoxStorageService} la gestión de cada tipo de archivo:

\begin{itemize}
    \item \texttt{CommunityRegistrationStorageService}: Documentos de registro de comunidades (formularios, actas constitutivas).
    \item \texttt{CommonExpenseReceiptStorageService}: Comprobantes de pago subidos por los residentes.
    \item \texttt{MarketplaceStorageService}: Imágenes de artículos publicados en el marketplace, optimizadas automáticamente por \texttt{ImageOptimizer} antes de la subida.
\end{itemize}

Los archivos se organizan en carpetas jerárquicas dentro de Box (por edificio, módulo y entidad), y se recuperan mediante el identificador de archivo almacenado en la base de datos. Esta arquitectura permite centralizar el almacenamiento en un solo proveedor, simplificando la gestión de permisos y el mantenimiento del sistema.

\subsection{Migraciones de Base de Datos}
El esquema de base de datos se gestiona mediante \textbf{Flyway}~\cite{flyway2025}, integrado directamente en el proceso de arranque de la aplicación a través del módulo de inyección de dependencias. Al iniciar el servidor, Flyway ejecuta automáticamente las migraciones pendientes, asegurando que el esquema esté siempre sincronizado con la versión del código.

El proyecto cuenta con 35~migraciones versionadas (V001 a V035) que documentan la evolución completa del modelo de datos: desde la creación del esquema de autenticación y roles, pasando por la incorporación de visitas, incidentes, votaciones, áreas comunes, finanzas, marketplace, chat, foro, encomiendas, tareas, personal, biblioteca, notificaciones, proveedores y órdenes de servicio, y las optimizaciones de rendimiento (índices y ajustes de esquema). El detalle completo de cada migración se presenta en el Anexo~\ref{anexo:diccionario}.

\section{Conclusión del Capítulo}
El sistema \textbf{DOMU} ha completado las siete fases de desarrollo planificadas, alcanzando un estado funcional completo con más de 130~endpoints REST, dos canales WebSocket, 35~migraciones de base de datos y una interfaz de usuario compuesta por 54~vistas. Los módulos de autenticación, gestión de comunidades, control de acceso, finanzas, incidentes, votaciones, reservas, chat, foro, marketplace, encomiendas, tareas, personal, biblioteca, notificaciones, proveedores y órdenes de servicio se encuentran operativos e integrados. Adicionalmente, la API se documenta de forma interactiva mediante Swagger~UI. El siguiente paso corresponde a las pruebas de aceptación, el despliegue en producción y la marcha blanca descritos en el capítulo siguiente.
